Transistor-level transient simulation repeatedly solves nonlinear modified nodal equations. A poor initial state increases device evaluations, Jacobian assemblies, and linear solves at each time point. Digital circuits expose a structure that generic warm-start methods miss: a small library of standard cells repeats across a large netlist, while shared nets couple their local states. We present \system, a learned warm-start method that follows this hierarchy. A type-shared proposer predicts a local correction for each cell instance from a three-step solver history, continuation state, and one-hop context. A conflict-aware aggregator combines overlapping proposals through pin roles, proposal disagreement, and the netlist graph. \system changes only the Newton initial state. The native \spice{} equations, convergence tests, and fallback path determine every accepted solution. Across \NumTestCircuits{} circuits in a licensed 40-nm library and the public ASAP7 kit, \system achieves \HeadlineSpeedup{} end-to-end speedup and reduces Newton iterations by \NewtonReduction{}. It covers \CoverageForty{} of validated 40-nm cells, preserves the baseline convergence rate within \ConvergenceDelta{}, and limits maximum waveform error to \WaveformAbsError{}. These results show that cell semantics provide a reusable unit for solver assistance across circuit sizes and process libraries.

